% !TeX program = pdflatex
% Build: bash build_paper.sh paper_en
% (pdflatex is run until the cross-references converge; two passes are not
% enough, three are needed.)
%
% This is the international-format counterpart of paper_ru.tex. It carries no
% GOST Р 7.0.7-2021 apparatus (article type, UDC index, bilingual captions,
% romanized reference list, "Information about the author" block); instead it
% follows the layout expected by international journals: single-language
% captions, a Declarations block, and one reference list with DOIs.
\documentclass[11pt,a4paper]{article}

% cmap MUST precede fontenc: it embeds ToUnicode tables in the PDF, without
% which the ff/fi/fl ligatures cannot be extracted -- the text is then neither
% searchable nor correctly indexed by Crossref, Google Scholar and the like.
\usepackage{cmap}
\usepackage[T1]{fontenc}
\usepackage[utf8]{inputenc}
\usepackage[english]{babel}
\usepackage{amsmath,amssymb,amsthm}
\usepackage{booktabs}
\usepackage{graphicx}
% Figures are generated by analysis.py; the English-labelled set lives in
% figs/en/, the Russian one in figs/.
\graphicspath{{figs/en/}{figs/}{./}}
\usepackage[margin=25mm]{geometry}
% xurl BEFORE hyperref: it allows long URLs to break anywhere. Without it,
% links such as t5k.org/mersenne/heuristic.html run into the right margin.
\usepackage{xurl}
\usepackage[colorlinks=true,allcolors=blue]{hyperref}
\emergencystretch=3em
\usepackage{xcolor}
\usepackage[expansion=false]{microtype}
\usepackage{caption}
\captionsetup{labelsep=period}

\theoremstyle{plain}
\newtheorem{proposition}{Proposition}
\newtheorem{corollary}{Corollary}
\theoremstyle{remark}
\newtheorem{remark}{Remark}

\newcommand{\Rb}{R_b}
\newcommand{\eg}{e^{\gamma}}
\newcommand{\emg}{e^{-\gamma}}
\DeclareMathOperator{\Gam}{Gamma}

\title{The observation scheme in the statistics of generalized\\
repunit primes: calibrating the pooled test of the\\
Lenstra--Pomerance--Wagstaff constant}
\author{Timofei S. Dokuchaev\thanks{Independent researcher, Irkutsk, Russia.
E-mail: \texttt{timofei93@yandex.ru}. ORCID:
\url{https://orcid.org/0009-0006-0510-5225}.}}
% Fixed date: \today would change on every rebuild.
\date{27 August 2026}

\begin{document}
\maketitle

\begin{abstract}
\noindent The Lenstra--Pomerance--Wagstaff (LPW) heuristic predicts
$\#\{n\le x:\Rb(n)\ \text{prime}\}\sim(\eg/\ln b)\ln x$, where
$\Rb(n)=(b^{n}-1)/(b-1)$. Since 1993 this universality has been tested through a
``normalized occurrence frequency'' obtained by regressing $\ln n_k$ on $k$.

First, the suspicion that right-censoring biases this frequency downwards is
\emph{false}: both the endpoint form $D/(N\ln b)$ and the OLS slope are exactly
unbiased (Proposition~\ref{prop:unbiased}); only the reciprocal is biased, and
that is Jensen's inequality rather than censoring. The genuine defect is
different: competing forms of the estimator disagree by up to $29.5\%$ on
identical data---$1.7$ times the median deviation from $\emg$---and the ``lucky
base'' label flips when the form is changed.

Second, an exact conditional likelihood replaces the regression and pools
$B=15$ bases into $M=218$ events, against $51$ for the Mersenne sequence alone.
Our main result concerns the pivot $\kappa S\sim\Gam(M,1)$: it is exact only
under inverse sampling (``observe until the $N$-th event''), whereas the data
arise from a fixed search frontier with a random number of finds. Under the
actual scheme a nominal $5\%$ test rejects in $16.2\%$ of cases, and the correct
pooled estimator is $(M-B)/S$, not $(M-1)/S$: the conditional likelihood costs
one event \emph{per base}. This yields $\hat\kappa=1.838$ with a $95\%$ interval
$[1.61,\,2.13]$ and $p=0.64$ against $\eg\approx1.781$; the reported
ten-percent excess over LPW is an artifact of the observation scheme.

Third, the choice between the two estimators is settled by how the data were
collected and not by the data themselves, and it costs $7\%$; hence our call to
publish the search frontiers $L_b$. Finally, the homogeneity test
($\Lambda=7.40$, $p=0.92$) resolves only a spread of $\kappa_b$ above
$\pm34\%$, and the pooled sample only deviations from $\eg$ above $21\%$.
\end{abstract}

\medskip
\noindent\textbf{Keywords:} repunit primes, generalized repunits, Mersenne
primes, Lenstra--Pomerance--Wagstaff conjecture, Euler--Mascheroni constant,
observation scheme, Type-I censoring, inverse sampling, Poisson point process,
conditional likelihood, sufficient statistic, estimator bias, statistical power,
computational number theory

\medskip
\noindent\textbf{Mathematics Subject Classification (2020):}
Primary 11Y11; Secondary 11N05, 62F03, 62B05, 11N80.

\section{Introduction}

For an integer base $b\ge2$ the generalized repunit is
\begin{equation}\label{eq:Rb}
\Rb(n)=\frac{b^{n}-1}{b-1},
\end{equation}
and $\Rb(n)$ can be prime only when $n$ is prime. For $b=2$ these are the
Mersenne numbers $M_p=2^{p}-1$. Building on the heuristics of
Wagstaff~\cite{wagstaff1983} and Pomerance~\cite{pomerance1981}, the LPW
conjecture asserts that
\begin{equation}\label{eq:LPW}
\pi_b(x)=\#\{n\le x:\ \Rb(n)\ \text{prime}\}\ \sim\ \frac{\eg}{\ln b}\,\ln x,
\end{equation}
where $\gamma$ is the Euler--Mascheroni constant; a detailed derivation of the
heuristic for $b=2$ is given in the PrimePages reference~\cite{t5k}. The base
enters~\eqref{eq:LPW} only through the factor $1/\ln b$: the constant
$\eg\approx1.7811$ is declared universal. Dubner~\cite{dubner1993} proposed
testing this universality empirically across bases, and
Bourdelais~\cite{bourdelais2009} has since maintained running estimates of the
normalized frequency, publishing them in the comments to OEIS
entries~\cite{oeis}. Related questions for the Mersenne sequence alone were
discussed recently in~\cite{wagstaff2025}. Close in spirit is
also~\cite{dominguez2026}, which looks for secondary structure in the Mersenne
exponents related to the factorization of $p-1$; methodologically it is built
like our~\S\ref{sec:struct}---stratified conditional likelihood plus a
significance test---but it tests for structure \emph{within} a single base,
whereas we test for it \emph{between} bases.

An interpretive folklore has grown up around these estimates. Base~13 is
described as ``lucky'' and base~7 as ``the unluckiest''; it is noted that the
normalized frequencies lie below the conjectured limit $\emg\approx0.5615$ and
converge to it slowly from below. The present work examines these claims, but
arrives at conclusions other than those one would naturally expect.

\subsection{Contributions}

\begin{enumerate}
\item\textbf{The estimator in common use is not biased.} The scheme in which the
fitting interval ends at the largest known index does indeed condition the
estimate on the last observation being an event, and the sufficient statistic is
distributed as $\Gam(N,\lambda_b)$ (Proposition~\ref{prop:gamma}). But this does
\emph{not} imply that the normalized frequency is biased: both the endpoint form
$D/(N\ln b)$ and the OLS slope are exactly unbiased
(Proposition~\ref{prop:unbiased}, Table~\ref{tab:mc}). Only the reciprocal
$\hat\kappa=1/\hat G$ is biased, by the factor $N/(N-1)$, and that is Jensen's
inequality, not censoring. Unbiased estimators of $\kappa$ and of $1/\kappa$
cannot be reciprocals of one another; any single ``correction'' that repairs one
of them damages the other (\S\ref{sec:notbias}).

\item\textbf{The genuine defect is instability with respect to the form of the
estimator.} Several numerically distinct variants are in circulation. On
identical data they disagree by as much as $29.5\%$ (base~18)---$1.7$ times the
median deviation of the estimates from $\emg$ ($17.6\%$), which is what gets
interpreted substantively. For two of the fifteen bases the choice of form
weighs more than their entire deviation from the LPW prediction. Under the
endpoint form base~18 is the third ``unluckiest'' of fifteen, and under OLS it is
the tenth, below the median. The label describes the choice of estimator, not a
property of the primes (\S\ref{sec:instability}).

\item\textbf{Pooling the bases, and what it costs.} The conditional likelihood
(\S\ref{sec:condlik}) does not require knowledge of the undocumented search
bounds $L_b$ and makes it possible to pool $B=15$ bases into $M=218$ events,
against $51$ for the Mersenne sequence alone. The exact null distribution
$\Gam(M,1)$ makes simulation unnecessary.

\item\textbf{Main result: the pivot is not calibrated to the actual observation
scheme.} The conditional likelihood is exact under the scheme ``observe until
the $N$-th event'', where $N$ is fixed by design. The data, however, arise from
a fixed frontier: the search covered $(t_0,\ln L_b]$ and the number of finds is
random. Under that scheme a nominal $5\%$ test rejects in $16.2\%$ of cases, and
the correct pooled estimator is $(M-B)/S$ (Proposition~\ref{prop:schemeB}). The
distinction between the schemes is itself classical, and the unbiasedness of
$(N-1)/D$ under scheme~A is a textbook fact to which we lay no claim; what is
new is that the frontier is \emph{undocumented}, so the exposure must be
reconstructed from the finds themselves, and the correction turns out to be
proportional to the number of bases (Remark~\ref{rem:priority}). This gives
$\hat\kappa=1.838$ with $p=0.64$: the excess over LPW disappears
(\S\ref{sec:scheme}).

\item\textbf{The data do not settle the choice of scheme.} The estimators
$(M-1)/S$ and $(M-B)/S$ are built on the same sufficient statistics $(N_b,D_b)$
and differ by $7\%$; which of them is right is determined not by the data but by
how the data were collected. An attempt to sidestep the choice by assembling the
pooled estimate from per-base pieces that are unbiased by construction fails
(\S\ref{sec:noescape}). The only radical remedy is to publish the frontiers
$L_b$.

\item\textbf{Lucky bases: what the test cannot do.} The likelihood-ratio test
gives $\Lambda=7.40$ on $14$ degrees of freedom, $p=0.92$. But its power is such
that it resolves only a spread of $\kappa_b$ of order $\pm34\%$
(\S\ref{sec:homog}). The correct conclusion is ``base-dependent frequencies are
not required by the data'', not ``they do not exist''.

\item\textbf{An explicit power calculation.} The pooled sample resolves
deviations from $\eg$ of at least $21.2\%$; the Mersenne sequence alone, at
least $49.8\%$ (\S\ref{sec:power}).
\end{enumerate}

Section~\ref{sec:sameerror} works through a narrower but instructive episode: a
naive regression of the Mersenne counting function produces an apparently
significant negative intercept, which turns out to be an artifact of
conditioning. We include it as a worked example of the same error in a different
guise.

\section{Model and data}\label{sec:data}

\subsection{The Poisson model}
Put $t=\ln n$. The heuristic~\eqref{eq:LPW} is equivalent to the statement that
the indices of generalized repunit primes to base~$b$ form a homogeneous Poisson
point process on the $t$-axis with constant intensity
\begin{equation}\label{eq:lam}
\lambda_b=\frac{\kappa}{\ln b},\qquad \text{LPW:}\ \ \kappa=\eg=1.781072\ldots
\end{equation}
Universality is the assertion that a single parameter $\kappa$ governs all
bases. The Poisson idealization here is a working assumption, not a theorem:
even for the ordinary primes the Poisson law in short intervals has been
established only conditionally, under the Hardy--Littlewood conjectures, and it
is known to break down as the scale grows~\cite{jha2026}. We adopt~\eqref{eq:lam}
as a model and treat $\kappa$ throughout as a quantity to be estimated; this
reduces the question ``is LPW true?'' to a one-parameter one and makes the bases
poolable.

\subsection{Admissible bases}\label{sec:bases}
Bases that are perfect powers are excluded: for $b=c^{m}$ with $m>1$ the number
$\Rb(n)$ admits an algebraic factorization and the process does not have the
assumed form. Bases $4,8,9,16$ are therefore omitted. Within the declared range
the selection does not depend on the results: we take \emph{all} bases
$2\le b\le20$ that are not perfect powers, which gives $B=15$; no base was
excluded or added on the strength of its estimate. Each of them has at least
eight known indices (the minimum is attained at $b=18$: eight terms of A133857).

We stress, however, that the upper limit $b\le20$ is itself a conventional round
number and not a consequence of the criterion ``at least eight indices'': that
criterion continues to hold beyond it (A127996 for $b=21$ has eight terms,
A127997 for $b=22$ has nine) and first fails at $b=23$ (A204940, six terms). To
make sure the limit does not look chosen after the fact, we tested its influence
directly. Extending the set to all bases $b\le26$ that are not perfect powers
(this adds $b=21,22,23,24,26$; $b=25=5^2$ is excluded algebraically) gives
$B=20$, $M=253$ and
\[
  \frac{M-1}{S}=1.9776,\qquad \frac{M-B}{S}=1.8285,
\]
that is, it shifts the main estimate~\eqref{eq:resB} by $0.5\%$---more than an
order of magnitude less than the price of the choice of observation scheme
(\S\ref{sec:scheme}). The homogeneity test on the extended set gives
$\Lambda=10.27$ on $19$ degrees of freedom, $p=0.95$. No conclusion of this work
depends on the limit $b\le20$; we retain it as pre-declared.

\subsection{The data and their status}
Table~\ref{tab:data} lists the sequences used, retrieved on 27 August 2026
(cross-check: Appendix~\ref{app:data}). Three caveats matter.

\paragraph{Probable primes.} For large indices most of the entries are probable
primes established by a Fermat or strong Fermat test after trial division. The
count is biased upwards only if a Fermat pseudoprime has been mistaken for a
prime, which for numbers of this size is practically unheard of. Sensitivity is
examined in~\S\ref{sec:robust}.

\paragraph{Verification frontiers.} For $b=2$ the GIMPS project~\cite{gimps}
publishes two frontiers: all exponents below $X_{\mathrm{ver}}=81\,648\,221$
have been tested twice, which gives exactly $50$ Mersenne primes, whereas all
exponents below $X_{\mathrm{ft}}=141\,308\,443$ have been tested at least once,
which gives $52$. The two largest exponents are provisional: unverified
candidates remain between them. The analysis of~\S\ref{sec:sameerror}, where the
frontier is used explicitly, is restricted to the $50$ verified exponents. The
pooled analysis is insensitive to this choice (\S\ref{sec:robust}).

Both frontiers are moving quantities, so we fix the moment of reading: the
PrimeNet Milestones report, 27 August 2026. No conclusion depends on the exact
value of $X_{\mathrm{ver}}$: it suffices that it lies between
$M_{50}=77\,232\,917$ and $M_{51}=82\,589\,933$, and this holds for every value
the report has taken during 2026.

\paragraph{The status of entry A000043.} The DATA section of that entry contains
$50$ terms and stops at $77\,232\,917$: OEIS regards the order of the subsequent
exponents as unproven until double-checking is complete. The exponents
$82\,589\,933$ and $136\,279\,841$ are known and documented by GIMPS (and are
mentioned in the comments to A000043), but formally they are not part of DATA.
We include them, because the statistic~\eqref{eq:pivot} uses only the position of
the last event and not its ordinal number; \S\ref{sec:robust} verifies that
restricting base~2 to the fifty double-checked exponents does not change the
conclusion.

\paragraph{Search bounds for $b>2$.} For the other bases no comparable frontier
is published. It is precisely this gap that is the subject of
\S\ref{sec:scheme}, and, we believe, the principal obstacle to a quantitative
test of LPW.

\begin{table}[t]
\centering
\caption{Sequences used. $N_b$ is the number of events strictly to the right of
$t_0=\ln n_{\min}$, that is, one less than the number of known prime indices:
the smallest index sets the origin, not an event. In the last row
$M=\sum_b N_b$ and $S=\sum_b D_b/\ln b$; the quantity $S$ is not the sum of the
$D_b$ column. Source: OEIS entries, retrieved 27 August 2026.}
\label{tab:data}
\small
\begin{tabular}{rlrrrrr}
\toprule
$b$ & OEIS & primes & $N_b$ & $n_{\min}$ & $n_{\max}$ & $D_b$\\
\midrule
2 & A000043 & 52 & 51 & 2 & 136\,279\,841 & 18.0371\\
3 & A028491 & 23 & 22 & 3 & 8\,530\,117 & 14.8605\\
5 & A004061 & 19 & 18 & 3 & 3\,300\,593 & 13.9110\\
6 & A004062 & 17 & 16 & 2 & 3\,360\,347 & 14.3344\\
7 & A004063 & 10 & 9 & 5 & 1\,264\,699 & 12.4409\\
10 & A004023 & 11 & 10 & 2 & 8\,177\,207 & 15.2237\\
11 & A005808 & 13 & 12 & 17 & 1\,868\,983 & 11.6077\\
12 & A004064 & 14 & 13 & 2 & 769\,543 & 12.8604\\
13 & A016054 & 13 & 12 & 5 & 1\,503\,503 & 12.6139\\
14 & A006032 & 11 & 10 & 3 & 1\,724\,417 & 13.2618\\
15 & A006033 & 10 & 9 & 3 & 639\,833 & 12.2704\\
17 & A006034 & 12 & 11 & 3 & 1\,990\,523 & 13.4053\\
18 & A133857 & 8 & 7 & 2 & 1\,270\,141 & 13.3615\\
19 & A006035 & 11 & 10 & 19 & 209\,359 & 9.3074\\
20 & A127995 & 9 & 8 & 3 & 984\,349 & 12.7011\\
\midrule
\multicolumn{3}{l}{total} & $M=218$ & & & $S=110.4423$\\
\bottomrule
\end{tabular}
\end{table}

\section{What is biased and what is not}\label{sec:notbias}

\subsection{The estimators in common use}
Dubner~\cite{dubner1993} and Bourdelais~\cite{bourdelais2009} characterize a
base~$b$ by a normalized occurrence frequency obtained by fitting a straight
line to $\ln n_k$ against $k$ and dividing by $\ln b$. By~\eqref{eq:lam} the
expected slope is $1/\lambda_b=\ln b/\kappa$, so the normalized frequency
estimates
\begin{equation}\label{eq:G}
G_b=\frac{1}{\kappa},\qquad\text{LPW:}\ \ G_b=\emg=0.5614595\ldots
\end{equation}
Several variants are in circulation---ordinary OLS with an intercept, regression
through the origin, and the endpoint form $\hat G_b=D_b/(N_b\ln b)$ with
$D_b=t_{N_b}-t_0$. All of them share one structural feature: the fitting
interval ends at the largest known index.

\subsection{The exact distribution}

\begin{proposition}[Conditional law of the right endpoint]\label{prop:gamma}
Let the events form a homogeneous Poisson process of intensity $\lambda$ on
$(t_0,\infty)$, and let $t_1<\dots<t_N$ be its first $N$ events after $t_0$.
Then $D:=t_N-t_0$ has the $\Gam(N,\lambda)$ distribution with density
$\lambda^{N}d^{N-1}e^{-\lambda d}/(N-1)!$. Conditionally on $D=d$, the preceding
$N-1$ events are distributed as the order statistics of $N-1$ independent
uniform variables on $(t_0,d)$.
\end{proposition}

\begin{proof}
The waiting times $t_1-t_0,\dots,t_N-t_{N-1}$ are independent
$\mathrm{Exp}(\lambda)$ variables, and a sum of $N$ of them is
$\Gam(N,\lambda)$. Conditional uniformity is the order-statistics property of
the Poisson process, applied to $(t_0,d)$.
\end{proof}

\begin{corollary}[Bias of the reciprocal]\label{cor:kappa}
The maximum-likelihood estimator of $\lambda$ based on $D$ alone is
$\hat\lambda=N/D$, and
\begin{equation}\label{eq:bias}
\mathbb E[\hat\lambda]=N\,\mathbb E[D^{-1}]=\frac{N}{N-1}\,\lambda,
\end{equation}
since $\mathbb E[D^{-1}]=\lambda/(N-1)$ for $D\sim\Gam(N,\lambda)$. The
estimator $\tilde\lambda=(N-1)/D$ is exactly unbiased.
\end{corollary}

\subsection{What does not follow from Corollary~\ref{cor:kappa}}

It is natural to conclude that, since $\hat\lambda$ is inflated by the factor
$N/(N-1)$, the frequency $\hat G_b$ must be deflated by the same factor. This is
false.

\begin{proposition}[Unbiasedness of the frequency]\label{prop:unbiased}
Under the hypotheses of Proposition~\ref{prop:gamma}:
\begin{enumerate}
\item[\rm(i)] the endpoint estimator is exactly unbiased:
$\mathbb E\!\left[\dfrac{D}{N\ln b}\right]=\dfrac{1}{\lambda\ln b}=G_b$;
\item[\rm(ii)] the OLS slope in the regression of $t_k$ on $k$ is exactly
unbiased: $\mathbb E[\hat\beta]=1/\lambda$;
\item[\rm(iii)] consequently the ``corrected'' quantity
$\tilde G_b=\frac{N}{N-1}\hat G_b$ is biased \emph{upwards}, by exactly the
factor $N/(N-1)$.
\end{enumerate}
\end{proposition}

\begin{proof}
(i) $\mathbb E[D]=N/\lambda$ for $D\sim\Gam(N,\lambda)$, whence the result at
once. (ii) The OLS slope is the linear statistic
$\hat\beta=\sum_k w_k t_k$ with weights
$w_k=(k-\bar k)/\sum_j(j-\bar k)^2$, for which $\sum_k w_k=0$ and
$\sum_k w_k k=1$. Since $\mathbb E[t_k]=t_0+k/\lambda$, we get
$\mathbb E[\hat\beta]=\sum_k w_k\,(t_0+k/\lambda)=1/\lambda$. Part (iii) is
immediate from (i).
\end{proof}

The substantive point is simple: unbiased estimators of $\kappa$ and of
$1/\kappa$ cannot be reciprocals of one another, because $x\mapsto1/x$ is
strictly convex. The unbiased pair is $(\hat G_b,\ \tilde\kappa_b)$ with
$\tilde\kappa_b=(N_b-1)\ln b/D_b$; but $\tilde\kappa_b\ne1/\hat G_b$. Any single
``censoring correction'' that repairs one of these quantities damages the other.
Table~\ref{tab:mc} confirms Proposition~\ref{prop:unbiased} by simulation.

\begin{table}[t]
\centering
\caption{Monte Carlo check of the propositions, $200\,000$ replications per row
(generator seed fixed, see \texttt{analysis.py}). The target value of every
column except the second is $1.0000$.}
\label{tab:mc}
\small
\begin{tabular}{rrrrrrr}
\toprule
$N$ & $\mathbb E[\hat\lambda]/\lambda$ & $N/(N-1)$ & $\mathbb E[\tilde\lambda]/\lambda$
& $\mathbb E[\hat G]\lambda$ & $\mathbb E[\tilde G]\lambda$ & $\mathbb E[\hat\beta]\lambda$\\
\midrule
7 & 1.1665 & 1.1667 & 0.9998 & 1.0002 & 1.1669 & 1.0002\\
10 & 1.1115 & 1.1111 & 1.0004 & 0.9998 & 1.1109 & 0.9997\\
16 & 1.0672 & 1.0667 & 1.0005 & 0.9994 & 1.0660 & 0.9994\\
22 & 1.0481 & 1.0476 & 1.0004 & 0.9996 & 1.0472 & 0.9997\\
51 & 1.0202 & 1.0200 & 1.0002 & 0.9999 & 1.0199 & 0.9998\\
\bottomrule
\end{tabular}
\end{table}

\begin{remark}
The negative result of this subsection deserves emphasis, because it runs
against intuition: the systematic downward deviation of the published
frequencies from $\emg$ is \emph{not} explained by the fitting scheme. Its
source is either a genuine excess of $\kappa$ over $\eg$ or, as
\S\ref{sec:scheme} shows, the observation scheme---but not the arithmetic of the
estimator.
\end{remark}

\subsection{Instability with respect to the form of the estimator}\label{sec:instability}

Unbiasedness does not mean agreement. Table~\ref{tab:bases} gives the endpoint
form $\hat G_b$ and the OLS estimate $\hat G_b^{\mathrm{OLS}}$ on identical
data. The discrepancies reach $29.5\%$ ($b=18$), $15.9\%$ ($b=13$) and $10.3\%$
($b=11$).

Compare these with what the literature interprets substantively---the deviations
of the estimates themselves from $\emg$. Across the fifteen bases the latter
have median $17.6\%$ (from $5.6\%$ at $b=20$ to $43.7\%$ at $b=19$). The largest
discrepancy between forms, $29.5\%$, exceeds that median by a factor of $1.7$.
On average the choice of form is of course cheaper: the median discrepancy is
$4.1\%$, about a quarter of the median deviation. But for two of the fifteen
bases ($b=18$ and $b=20$) the choice of form weighs more than their entire
deviation from the LPW prediction---and one of them, base~18, is precisely a base
the folklore discusses.

The consequence for the folklore is direct. Under the endpoint form base~18 has
$\hat G_{18}=0.6604$, the third largest of fifteen (after $b=7$ with $0.7104$
and $b=10$ with $0.6612$), so it belongs firmly among the three ``unluckiest'';
under OLS, $\hat G_{18}^{\mathrm{OLS}}=0.4656$, tenth of fifteen and below the
median $0.4904$, so it has become ``lucky''. The label flips without a single
new prime, purely from a change in the form of the estimator: the base moves
from third place to tenth.

To this is added a second, technical difficulty with OLS: its standard errors
are wrong in this problem. The residuals are deviations of order statistics and
are strongly autocorrelated; for $b=2$ we obtain $\rho_1=0.825$, and the
Newey--West covariance estimator with eight lags raises the standard error of
the intercept from $0.447$ to $0.732$ (\S\ref{sec:sameerror}). Published point
values come with no intervals at all, and this---not bias---is the principal
methodological objection to current practice.

\subsection{The conditional likelihood}\label{sec:condlik}

Proposition~\ref{prop:gamma} yields more than an analysis of bias. Setting
$t_0=\ln n_{\min}$ and taking $(N_b,D_b)$ as the data for base~$b$, we obtain
the log-likelihood
\begin{equation}\label{eq:loglik}
\ell(\kappa)=\sum_b\Bigl[N_b\ln\frac{\kappa}{\ln b}-\frac{\kappa}{\ln b}D_b\Bigr]+\mathrm{const},
\end{equation}
whose distribution does not depend on the search bound $L_b$. This matters in
practice: $L_b$ is well documented only for $b=2$. Differentiating
\eqref{eq:loglik} and putting $S=\sum_b D_b/\ln b$, we get
\begin{equation}\label{eq:mle}
\hat\kappa=\frac{M}{S},\qquad \tilde\kappa=\frac{M-1}{S},\qquad M=\sum_b N_b .
\end{equation}
Since $\kappa S=\sum_b\lambda_bD_b$ is a sum of independent $\Gam(N_b,1)$
variables,
\begin{equation}\label{eq:pivot}
\kappa S\sim\Gam(M,1),
\end{equation}
an exact pivot; confidence intervals and $p$-values follow in closed form.
\emph{It is precisely this statement that requires revision}---see
\S\ref{sec:scheme}.

\begin{remark}[Independence across bases]
Pooling presupposes that the processes for different bases are independent.
Strictly this fails at small indices: seventeen index values occur in more than
one base, and $n=3$ occurs in nine of them, since $\Rb(3)=b^2+b+1$ is prime for
many~$b$. All such coincidences lie at $n\le317$. More important than coinciding
indices (the numbers $\Rb(n)$ for different $b$ are distinct) is the possible
sharing of algebraic divisors; both affect only small~$n$, whereas $D_b$ is
determined by the upper end of the range. Section~\ref{sec:robust} confirms that
discarding all indices below~$10^3$ leaves the conclusion unchanged.
\end{remark}

\section{Two observation schemes}\label{sec:scheme}

\subsection{The distinction}

The pivot~\eqref{eq:pivot} was derived on the assumption that $N_b$ is fixed by
the observation plan: we watch the process until $N_b$ events have accumulated,
then stop. Call this \emph{scheme~A} (inverse sampling, Type-II censoring).

The data arose differently. For each base the computational search covered the
whole interval $(t_0,\ln L_b]$, and every event falling inside it was found; the
number of finds $N_b$ is a random variable, while the frontier $L_b$ is fixed.
Call this \emph{scheme~B}.

\begin{remark}[What here is classical and what is not]\label{rem:priority}
The distinction between the schemes is itself classical and goes back to
reliability theory, where scheme~A is called a \emph{failure-truncated} and
scheme~B a \emph{time-truncated} life test~\cite{epstein1954trunc,%
epsteinsobel1954}. The estimator $(N-1)/D$, unbiased under scheme~A
(Corollary~\ref{cor:kappa}), is a textbook fact there, and we lay no claim to
it. It is also known that under scheme~B no exactly unbiased estimator of the
\emph{mean} $1/\lambda$ exists at all~\cite{bartholomew1963,basu1964}; that is a
warning about how delicate the question of unbiasedness becomes under Type-I
censoring. We stress, however, that the remainder in~\eqref{eq:EBd} has nothing
to do with that theorem: by Remark~\ref{rem:atom} it is the expected number of
bases yielding no find at all, and its origin is quite different.

What we regard as new is something else. In the classical setting the moment at
which the test is stopped is \emph{known}: the time $T$ is set by the
experimenter. Here the frontier $L_b$ is undocumented for every base except
$b=2$, so the full exposure cannot be read off a protocol and must instead be
reconstructed---via the memorylessness of the process---from the finds
themselves. That is where a correction \emph{per base}, rather than one for the
sample as a whole, comes from: there are as many unknown frontiers as there are
bases. A survey of the biases generated by the edge of the observation window,
though in a nonparametric setting and for the shape of the interevent
distribution rather than for the intensity, is given in~\cite{kivela2015}.
\end{remark}

The distinction is not cosmetic. Under scheme~B, conditionally on $N_b=n$, the
positions of the events are $n$ independent uniform variables on
$(t_0,\ln L_b]$, so the conditional distribution of $D_b$ depends on $L_b$ and
\emph{not} on $\lambda_b$: with $N_b$ fixed, $D_b$ is informative about the
frontier, not about the intensity. Under scheme~B the
likelihood~\eqref{eq:loglik} is not a likelihood.

Note also that \S\ref{sec:sameerror} of the present work, where the frontier for
$b=2$ is known and used explicitly, rests on scheme~B. One cannot apply both
schemes in a single paper; below we adopt scheme~B as matching the way the data
were obtained, and calibrate the scheme~A results to it.

\subsection{The correct pooled estimator under scheme B}

The bias has a closed form, and an instructive one.

\begin{proposition}\label{prop:schemeB}
Suppose that under scheme~B the search for base~$b$ covered $(t_0,\ln L_b]$, let
$W_b=\ln L_b-t_0$ be the width of the window, and let
$\delta_b=\ln L_b-t_{N_b}$ be the distance from the frontier to the last event
found (with $\delta_b=W_b$ when $N_b=0$). Then
\begin{equation}\label{eq:delta}
\Pr(\delta_b>x)=e^{-\lambda_b x}\ \ (0\le x<W_b),\qquad
\Pr(\delta_b=W_b)=e^{-\lambda_b W_b},
\end{equation}
that is, $\delta_b$ is an exponential variable with parameter $\lambda_b$
truncated at $W_b$, and
\begin{equation}\label{eq:Edelta}
\mathbb E[\delta_b]=\frac{1-e^{-\lambda_b W_b}}{\lambda_b}.
\end{equation}
The full exposure is
$S_{\mathrm{full}}=S+\sum_b\delta_b/\ln b=\sum_b W_b/\ln b$, with
$\mathbb E[M]=\kappa S_{\mathrm{full}}$ identically, and
\begin{equation}\label{eq:EBd}
\mathbb E\Bigl[\kappa\sum_b\delta_b/\ln b\Bigr]=B-\sum_b e^{-\lambda_b W_b}.
\end{equation}
Equating $\kappa S_{\mathrm{full}}$ to $M$ and dropping the exponentially small
remainder gives the method-of-moments estimator
\begin{equation}\label{eq:MBS}
\hat\kappa_{\mathrm B}=\frac{M-B}{S}.
\end{equation}
\end{proposition}

\begin{proof}
For $x<W_b$ the event $\{\delta_b>x\}$ means that the interval
$(\ln L_b-x,\ln L_b]$ of length $x$ contains no event of the process, which has
probability $e^{-\lambda_b x}$. As $x\to W_b^-$ what remains is the case ``no
events at all'', which gives the atom in~\eqref{eq:delta}. Integrating the tail,
$\mathbb E[\delta_b]=\int_0^{W_b}e^{-\lambda_b x}\,dx=(1-e^{-\lambda_b W_b})/\lambda_b$,
which is~\eqref{eq:Edelta}. Next, $S_{\mathrm{full}}$ is not random: by
construction $D_b+\delta_b=W_b$, so
$\kappa S_{\mathrm{full}}=\sum_b\lambda_b W_b=\mathbb E[M]$---this is exactly the
moment equation. Finally
$\kappa\,\mathbb E[\delta_b]/\ln b=\lambda_b\mathbb E[\delta_b]=1-e^{-\lambda_b W_b}$,
and summation over $b$ gives~\eqref{eq:EBd}.
\end{proof}

\begin{remark}\label{rem:atom}
The remainder $\sum_b e^{-\lambda_b W_b}$ dropped in~\eqref{eq:EBd} is the
expected number of bases for which the search produced no find at all. In our
data every base has finds and $\lambda_b W_b\approx N_b\ge7$, so the remainder
does not exceed $\sum_b e^{-N_b}<3\cdot10^{-3}$ against $B=15$---a relative
correction of order $0.02\%$. That is what the first row of
Table~\ref{tab:frontier} exhibits. We stress that \emph{without} the truncation
the statement $\delta_b\sim\mathrm{Exp}(\lambda_b)$ would be false: the process
has no events to the left of $t_0$, so $\delta_b\le W_b$ by construction, and an
appeal to memorylessness alone is not enough here.
\end{remark}

In other words, the conditional likelihood costs information equivalent to one
event \emph{per base}---whereas formula~\eqref{eq:mle} subtracts one event
\emph{in total}. With $B=15$ the difference is fourteen events, that is, $7\%$:
\begin{equation}
\frac{M-1}{S}=1.9648\qquad\text{versus}\qquad\frac{M-B}{S}=1.8381 .
\end{equation}
Direct simulation of scheme~B confirms this: estimator~\eqref{eq:MBS} has a bias
of $+0.02\%$, whereas $(M-1)/S$ has a bias of $+6.9\%$
(Table~\ref{tab:calib}).

It is essential that the conclusion of Proposition~\ref{prop:schemeB} rests
\emph{only} on identity~\eqref{eq:Edelta}, which holds for \emph{any} position of
the frontier: only the window width $W_b$ enters it, and that cancels together
with the exponentially small remainder (Remark~\ref{rem:atom}). Estimator
\eqref{eq:MBS} therefore contains no free parameter and does not require knowing
how far the frontier has advanced. Table~\ref{tab:frontier} confirms this by
simulation under seven different rules for placing the frontier, including
random ones and ones differing between bases.

\begin{table}[t]
\centering
\caption{Unbiasedness of the two pooled estimators under different frontier
placements. True value $\eg=1.7811$; $15\,000$ replications per row. The rules
are stated through the window width $W_b=\ln L_b-t_0$: ``$\mathbb E[\#]$ =
observed'' is $W_b=N_b/\lambda_b$; ``$+s$ mean intervals'' is
$W_b=(N_b+s)/\lambda_b$; ``$c$ times further in $t$'' is $W_b=c\,D_b$ with the
observed $D_b$; ``random'' is $W_b=(N_b/\lambda_b)\,U$ with independent
$U\sim\mathrm{Unif}(0.7,\,1.8)$, drawn separately for each base and each
replication.}
\label{tab:frontier}
\small
\begin{tabular}{lrrrr}
\toprule
Rule for the frontier & $\mathbb E[(M{-}B)/S]$ & dev. & $\mathbb E[(M{-}1)/S]$ & dev.\\
\midrule
$\mathbb E[\#\text{events}]$ = observed & 1.7814 & $+0.02\%$ & 1.9043 & $+6.92\%$\\
frontier $+0.5$ mean intervals & 1.7816 & $+0.03\%$ & 1.9001 & $+6.68\%$\\
frontier $+2$ mean intervals & 1.7820 & $+0.05\%$ & 1.8891 & $+6.07\%$\\
frontier $+5$ mean intervals & 1.7820 & $+0.05\%$ & 1.8717 & $+5.09\%$\\
frontier $1.5$ times further in $t$ & 1.7825 & $+0.08\%$ & 1.8716 & $+5.08\%$\\
frontier twice as far in $t$ & 1.7816 & $+0.03\%$ & 1.8475 & $+3.73\%$\\
random frontiers, differing by base & 1.7812 & $+0.01\%$ & 1.8788 & $+5.49\%$\\
\bottomrule
\end{tabular}
\end{table}

\paragraph{The only assumption.} All that is required is the following.

\medskip
\noindent\textbf{(A1)} \emph{The position of the frontier $L_b$ was not chosen in
the light of the search results: the computational budget, not the finds,
determines how far the sieving has been carried.}

\medskip
\noindent In the language of survival analysis, (A1) is the requirement that the
censoring be \emph{non-informative}: the stopping time does not depend on what
the process has produced. A violation of (A1) is exactly informative censoring
and, as usual in such problems, biases the estimate in the direction dictated by
the stopping rule; the size of the bias is written out below.

\medskip
\noindent If (A1) fails in the extreme form---the search for each base being
stopped immediately after a find---then $\delta_b=0$ and we are back at
scheme~A. The general form of the estimator when $\delta_b=m/\lambda_b$ is
$\hat\kappa=(M-Bm)/S$, so the whole range generated by the uncertainty in (A1)
is the segment between $1.838$ ($m=1$, frontier independent of the finds) and
$1.974$ ($m=0$, search stopped at a find).\footnote{%
For $m=0$ the moment equation gives exactly $M/S=1.9739$, whereas the exactly
unbiased scheme~A estimator of Corollary~\ref{cor:kappa} is $(M-1)/S=1.9648$
(the difference is one event out of $218$, that is, $0.5\%$; this is the usual
gap between a moment estimator and one corrected for Jensen's inequality). In
\S\ref{sec:pooled} and everywhere below we report $(M-1)/S$; here, for
homogeneity within the family $(M-Bm)/S$, the limit of the moment form is
written.}
Actual practice in distributed computing is clearly closer to the first: sieving
proceeds over contiguous intervals and finds are announced as they occur. It is
essential that both ends of the segment are compatible with LPW
(\S\ref{sec:pooled}), so the conclusion of this work does not depend on (A1);
what depends on (A1) is only the size of the residual gap.

\begin{table}[t]
\centering
\caption{Calibration of the pivot~\eqref{eq:pivot} under scheme~B. ``Frontier
shift'' is how many mean intervals beyond the expected last event $\ln L_b$ is
placed. $20\,000$ replications per row; Monte Carlo accuracy about $\pm0.005$
(binomial at the observed level $0.16$; the spread across five generator seeds
is $0.162$--$0.169$).}
\label{tab:calib}
\small
\begin{tabular}{lrr}
\toprule
Scheme & type~I error (nominal $0.05$) & $\mathbb E[\tilde\kappa]/\kappa=c$\\
\midrule
B, shift $0$ & 0.1616 & 1.0688\\
B, shift $0.5$ & 0.1600 & 1.0664\\
B, shift $1$ & 0.1558 & 1.0641\\
B, shift $2$ & 0.1491 & 1.0595\\
A (control) & 0.0504 & 0.9997\\
\bottomrule
\end{tabular}
\end{table}

Two conclusions follow from Table~\ref{tab:calib}. First, under scheme~B a
nominal $5\%$ test rejects a true hypothesis in $16.2\%$ of cases: the claimed
precision of the intervals is overstated roughly threefold. Second, under
scheme~B estimator~\eqref{eq:MBS} is practically unbiased, so the issue is not
an irremovable distortion but a wrongly chosen scheme.

The only direct check available is for $b=2$: there
$\delta_2=\ln(X_{\mathrm{ver}}/M_{50})=0.056$ against the mean interval
$\ln 2/\eg=0.389$ predicted by LPW. This is a single realization of an
exponential variable with unit mean, and such a value is not unusual:
$\Pr(\delta\lambda<0.143)=0.13$. (Taking the estimated rather than the predicted
mean interval, $1/\hat\lambda=0.358$, gives $\delta\lambda=0.155$ and
$\Pr=0.14$; the conclusion is the same.) Note that the unbiasedness of
estimator~\eqref{eq:MBS} does not depend on this realization---it averages over
all fifteen bases---and that the variability of $\delta_b$ itself is accounted
for in the confidence interval~\eqref{eq:resB}, obtained by inverting a
simulation of scheme~B.

\subsection{The data do not choose the scheme for us}\label{sec:noescape}

We stress that $(M-1)/S$ and $(M-B)/S$ are functions of the same sufficient
statistics $(N_b,D_b)$. No processing of the available data will say which of
them is right: the choice is determined by how the data were collected, and its
price is $7\%$---twice the entire residual gap with LPW.

The natural attempt to sidestep the choice is to build the pooled estimate out
of pieces that are unbiased by construction. The per-base estimator
$\tilde\kappa_b=(N_b-1)\ln b/D_b$ is unbiased under scheme~A
(Corollary~\ref{cor:kappa}) and unbiased under scheme~B up to an exponentially
small term. Indeed, for $N_b=n\ge2$ the ratio $D_b/W_b$ is distributed as
$\mathrm{Beta}(n,1)$, and the identity
$\mathbb E[1/\mathrm{Beta}(\alpha,\beta)]=(\alpha+\beta-1)/(\alpha-1)$ gives
$\mathbb E[\tilde\lambda_b\mid N_b=n]=n/W_b$; for $n\in\{0,1\}$ the estimator is
zero. Averaging over $n\sim\mathrm{Pois}(\lambda_bW_b)$ therefore gives not
$\lambda_b$ but
\[
  \mathbb E[\tilde\lambda_b]=\lambda_b-\frac{\Pr(N_b=1)}{W_b}
  =\lambda_b\Bigl(1-e^{-\lambda_bW_b}\Bigr),
\]
the same exponentially small deficit as in~\eqref{eq:EBd}: for
$\lambda_bW_b\approx N_b\ge7$ it does not exceed $0.1\%$. Simulation confirms
this: for each of the fifteen bases the mean of $\tilde\kappa_b$ lies within
$1.774$--$1.790$ against the true value $1.781$ (the estimator has infinite
variance when $N_b=2$, so the sample mean converges slowly; the residual spread
across bases is Monte Carlo noise).

Nevertheless, the weighted mean
$$\kappa_\ast=\sum\nolimits_b(N_b-2)\tilde\kappa_b\Big/\sum\nolimits_b(N_b-2)$$
with inverse-variance weights has a bias of $+0.14\%$ under scheme~A but
$+7.51\%$ under scheme~B, that is, essentially the same as $(M-1)/S$. The reason
is that the weights---the event counts themselves---are random under scheme~B
and positively correlated with the estimates; unbiasedness of each summand does
not carry over to a ratio of sums.

From this we draw a narrow but, we think, important conclusion: \emph{the gain
in precision from pooling the bases is inseparable from an assumption about the
observation scheme}. The only radical remedy is to document the search frontiers
$L_b$, as GIMPS does for $b=2$. We address this request to the participants of
the relevant search projects.

\section{Results}\label{sec:results}

\subsection{Per-base estimates}

\begin{table}[t]
\centering
\caption{Per-base estimates. $\hat G_b=D_b/(N_b\ln b)$ is the endpoint form;
$\hat G_b^{\mathrm{OLS}}$ is the OLS slope of $\ln n_k$ on $k$ divided by
$\ln b$; both are unbiased (Proposition~\ref{prop:unbiased}).
$\tilde\kappa_b=(N_b-1)\ln b/D_b$ is the unbiased estimator of the constant,
with an exact interval from $\Gam(N_b,1)$. The last column is the probability
integral transform of $\Gam(N_b,1)$ applied to $\eg D_b/\ln b$. LPW predicts
$G=\emg=0.5615$ and $\kappa=\eg=1.7811$.}
\label{tab:bases}
\small
\begin{tabular}{rrrrrrrlr}
\toprule
$b$ & $N_b$ & $\hat G_b$ & $\hat G_b^{\mathrm{OLS}}$ & diff., \% &
$\hat\kappa_b$ & $\tilde\kappa_b$ & 95\% CI & PIT\\
\midrule
2 & 51 & 0.5102 & 0.5309 & $+4.1$ & 1.960 & 1.921 & [1.459, 2.533] & 0.266\\
3 & 22 & 0.6148 & 0.6081 & $-1.1$ & 1.626 & 1.552 & [1.019, 2.373] & 0.693\\
5 & 18 & 0.4802 & 0.4666 & $-2.8$ & 2.083 & 1.967 & [1.234, 3.149] & 0.285\\
6 & 16 & 0.5000 & 0.4904 & $-1.9$ & 2.000 & 1.875 & [1.143, 3.092] & 0.355\\
7 & 9 & 0.7104 & 0.7371 & $+3.8$ & 1.408 & 1.251 & [0.644, 2.466] & 0.801\\
10 & 10 & 0.6612 & 0.6746 & $+2.0$ & 1.512 & 1.361 & [0.725, 2.584] & 0.737\\
11 & 12 & 0.4034 & 0.3617 & $-10.3$ & 2.479 & 2.272 & [1.281, 4.066] & 0.162\\
12 & 13 & 0.3981 & 0.3796 & $-4.7$ & 2.512 & 2.319 & [1.337, 4.050] & 0.141\\
13 & 12 & 0.4098 & 0.3448 & $-15.9$ & 2.440 & 2.237 & [1.261, 4.002] & 0.174\\
14 & 10 & 0.5025 & 0.5386 & $+7.2$ & 1.990 & 1.791 & [0.954, 3.400] & 0.406\\
15 & 9 & 0.5035 & 0.5054 & $+0.4$ & 1.986 & 1.766 & [0.908, 3.479] & 0.417\\
17 & 11 & 0.4301 & 0.4353 & $+1.2$ & 2.325 & 2.114 & [1.161, 3.887] & 0.229\\
18 & 7 & 0.6604 & 0.4656 & $-29.5$ & 1.514 & 1.298 & [0.609, 2.825] & 0.714\\
19 & 10 & 0.3161 & 0.2971 & $-6.0$ & 3.164 & 2.847 & [1.517, 5.405] & 0.061\\
20 & 8 & 0.5300 & 0.5713 & $+7.8$ & 1.887 & 1.651 & [0.815, 3.402] & 0.483\\
\bottomrule
\end{tabular}
\end{table}

Eleven of the fifteen bases lie below $\emg$---the same eleven under either form
of the estimator, so the sign of the deviation is robust to the choice of form,
unlike the ordering of the bases among themselves. All fifteen PIT quantiles lie
inside the central $95\%$ region; the extremes are $b=19$ at quantile $0.061$
and $b=7$ at $0.801$. Base~13, the ``lucky'' one, sits at $0.174$, which is in
no way remarkable for one base out of fifteen.

\subsection{The pooled estimate}\label{sec:pooled}

Under scheme~A, formulas~\eqref{eq:mle}--\eqref{eq:pivot} give
\begin{equation}\label{eq:resA}
\tilde\kappa=1.9648,\qquad 95\%\ \text{CI}=[1.7205,\,2.2444],\qquad p=0.142 .
\end{equation}
The point estimate is $10\%$ above $\eg$, and it is precisely this excess that
used to be interpreted as the trace of a factor omitted by the heuristic.

Estimator~\eqref{eq:MBS} for scheme~B gives (the interval and $p$-value are
obtained by inverting a simulation of scheme~B: for a grid of values of $\kappa$,
scheme~B is simulated with the observed frontier placement, and one seeks those
$\kappa$ at which the observed statistic $(M-B)/S$ falls on the $0.025$ and
$0.975$ quantiles of the simulated distribution; $20\,000$ replications per
point)
\begin{equation}\label{eq:resB}
\hat\kappa_{\mathrm B}=\frac{M-B}{S}=1.8381,\qquad
95\%\ \text{CI}=[1.613,\,2.130],\qquad p=0.639 .
\end{equation}
The excess falls from $10.3\%$ to $3.2\%$, and the $p$-value rises from $0.14$
to $0.64$. We hold that scheme~B is closer to the way the data were obtained,
and therefore regard~\eqref{eq:resB} as the main result and~\eqref{eq:resA} as an
upper bound: the latter corresponds to the assumption $m=0$, that is, to
stopping the search exactly at a find. The truth lies between them, and without
publication of the $L_b$ this interval cannot be narrowed. Both estimates are
consistent with LPW; they differ in whether the data call for an explanation of
a residual excess. We conclude that they do not. The overall picture is shown in
Figure~\ref{fig:kappa}: the confidence interval of every base covers both $\eg$
and both pooled estimates.

\begin{figure}[t]
\centering
\includegraphics[width=.85\textwidth]{fig_kappa.pdf}
\caption{Per-base estimates $\tilde\kappa_b$ with exact $95\%$ confidence
intervals from $\Gam(N_b,1)$; marker area is proportional to $N_b$. Dashed line:
the LPW value $\eg$; blue dash-dotted: the pooled estimate under scheme~A; green
with a hatched band: the estimate calibrated to scheme~B and its interval. Each
base separately is compatible with a universal constant.}
\label{fig:kappa}
\end{figure}

\subsection{Homogeneity, and its power}\label{sec:homog}

Let $H_0$ be the model with a single $\kappa$ and $H_1$ the model with free
$\kappa_b$. From~\eqref{eq:loglik}, since
$\sum_b\hat\kappa_bD_b/\ln b=M=\hat\kappa S$, the likelihood-ratio statistic
reduces to
\begin{equation}
\Lambda=2\sum_b N_b\ln\frac{\hat\kappa_b}{\hat\kappa}\ \overset{H_0}{\sim}\ \chi^2_{B-1}.
\end{equation}
We obtain $\Lambda=7.40$ on $14$ degrees of freedom, $p=0.92$.

But a non-rejection is informative exactly to the extent that the test is
powerful, and here one must apply to oneself the same requirement one applies to
the literature. We simulated alternatives in which the $\kappa_b$ are lognormally
scattered about $\eg$ with coefficient of variation $cv$:

\begin{center}
\small
\begin{tabular}{lrrrrrrrrr}
\toprule
$cv$ & 0.05 & 0.10 & 0.15 & 0.20 & 0.25 & 0.30 & 0.35 & 0.40 & 0.50\\
power & 0.061 & 0.120 & 0.212 & 0.363 & 0.535 & 0.693 & 0.817 & 0.893 & 0.965\\
\bottomrule
\end{tabular}
\end{center}

The test attains $80\%$ power only at $cv\approx0.34$ ($5000$ replications per
grid point, power accurate to about $\pm0.01$). Hence $p=0.92$ means: the data
do not require base-dependent frequencies, but neither do they exclude a spread
of $\kappa_b$ as large as a third of the value of the constant. The statement
``lucky bases do not exist'' does not follow from the available data; what
follows is the weaker statement ``their existence is not supported by the data,
and differences below $\pm34\%$ are in principle indistinguishable''.

\subsection{Robustness}\label{sec:robust}

\paragraph{Lower truncation.} We varied the lower truncation point in order to
probe the asymptotic regime (the heuristic~\eqref{eq:LPW} is asymptotic, whereas
the constant-intensity model starts at $n=2$). The results are in
Table~\ref{tab:trunc}. The scheme~A estimate grows from $1.965$ to $2.099$---not
strictly monotonically (at $n>10^{2}$ it retreats from $2.019$ to $2.012$) but
systematically: the growth is $6.8\%$, and all five values come from nested
subsamples of the same data. This is a trend, not noise. The scheme~B estimate
shows no trend and fluctuates about $\eg$ within $1.79$--$1.87$ (under the
weaker dropout criterion $N_b\ge1$, within $1.76$--$1.87$, still without a
trend). The reason is visible from~\eqref{eq:MBS}: truncation reduces $M$ but
not $B$, whereas scheme~A subtracts one regardless of the sample size. We regard
this as independent confirmation of the diagnosis of \S\ref{sec:scheme}: the
trend was a property of the estimator, not of the data. Both curves are shown in
Figure~\ref{fig:trunc}.

\begin{table}[t]
\centering
\caption{Robustness to lower truncation. A base is retained if at least
\emph{two} events survive the truncation ($N_b\ge2$); at $n>10^4$ base~19 no
longer meets this criterion, whence $B=14$. The criterion matters: under the
weaker requirement $N_b\ge1$ base~19 is retained and the last row becomes
$B=15$, $M=90$, $(M-1)/S=2.0848$, $(M-B)/S=1.7569$, $p=0.128$.}
\label{tab:trunc}
\small
\begin{tabular}{lrrrrr}
\toprule
truncation & $B$ & $M$ & $(M-1)/S$ & $(M-B)/S$ & $p$ (scheme A)\\
\midrule
none & 15 & 218 & 1.9648 & 1.8381 & 0.142\\
$n>10$ & 15 & 192 & 2.0192 & 1.8712 & 0.080\\
$n>10^2$ & 15 & 156 & 2.0116 & 1.8299 & 0.124\\
$n>10^3$ & 15 & 125 & 2.0551 & 1.8231 & 0.106\\
$n>10^4$ & 14 & 89 & 2.0990 & 1.7889 & 0.116\\
\bottomrule
\end{tabular}
\end{table}

\begin{figure}[t]
\centering
\includegraphics[width=.75\textwidth]{fig_truncation.pdf}
\caption{The pooled estimate under lower truncation: $(M-1)/S$ (scheme~A) and
$(M-B)/S$ (scheme~B). The trend of the scheme~A estimate disappears under
scheme~B.}
\label{fig:trunc}
\end{figure}

\paragraph{Leave-one-base-out.} A leave-one-base-out analysis gives
$\tilde\kappa$ within $[1.930,\,2.012]$ and $p$ between $0.086$ and $0.237$; no
single base determines the conclusion, and the most influential one, $b=3$,
shifts $p$ only to $0.086$ when removed.

\paragraph{Status of the data.} Restricting base~2 to the fifty double-checked
exponents moves the result from $\tilde\kappa=1.965$, $p=0.142$ to
$\tilde\kappa=1.961$, $p=0.151$. Excluding all indices $n>10^6$, where the
entries have PRP status, gives $\tilde\kappa=1.951$ with $M=182$. In the worst
case, when the last known PRP of some base turns out to be composite, the
largest shift comes from $b=17$: $\tilde\kappa=1.977$. All variants lie inside
interval~\eqref{eq:resA}, so possible pseudoprimes do not affect the conclusion.

\section{Statistical power}\label{sec:power}

By~\eqref{eq:pivot} the power against a true constant $\kappa=r\eg$ at level
$\alpha$ is available in closed form (Table~\ref{tab:power}).

\begin{table}[t]
\centering
\caption{Power at level $\alpha=0.05$ for rejecting $\kappa=\eg$ when the true
constant exceeds it by the stated factor.}
\label{tab:power}
\small
\begin{tabular}{lrr}
\toprule
$\kappa/\eg$ & Mersenne only ($M=51$) & pooled ($M=218$)\\
\midrule
1.05 & 0.060 & 0.106\\
1.10 & 0.095 & 0.277\\
1.15 & 0.152 & 0.523\\
1.20 & 0.229 & 0.755\\
1.30 & 0.427 & 0.972\\
1.50 & 0.802 & 1.000\\
\bottomrule
\end{tabular}
\end{table}

At $80\%$ power the Mersenne sequence alone detects only deviations of $49.8\%$
and above; pooling fifteen bases lowers that threshold to $21.2\%$. Attaining
$80\%$ power against a $10\%$ deviation would require $M\approx875$ events, and
against $5\%$, $M\approx3319$.

These thresholds are computationally out of reach, and it is worth saying by how
much. Since $M$ grows logarithmically in the search frontier, doubling the
frontiers \emph{of all fifteen bases} yields about $9.6$ additional events for
the entire pooled sample. To gather the missing $657$ events one needs
$657/9.6\approx68.2$ doublings, that is, a growth of the frontiers by a factor of
$2^{68.2}\approx3\cdot10^{20}$. This is not a matter of decades but a magnitude
lying beyond the reach of any conceivable computation. Both power curves, as
functions of the true deviation, are shown in Figure~\ref{fig:power}.

\begin{figure}[t]
\centering
\includegraphics[width=.72\textwidth]{fig_power.pdf}
\caption{Power of rejecting $\kappa=\eg$ at level $\alpha=0.05$. The horizontal
line marks $80\%$ power.}
\label{fig:power}
\end{figure}

Hence a reading rule for this literature. A report that the data agree with LPW
to within a few percent is not evidence for LPW at the level of a few percent;
it is compatible with a constant that is wrong by a factor of one and a half.
The pooled analysis does, however, establish a genuine bound: the universal
constant, if it exists, lies roughly in $[1.61,\,2.13]$, and no base deviates
from it detectably.

\section{The same error in a different guise}\label{sec:sameerror}

The Mersenne counting function invites a related error, which we present as a
worked example. Regressing $\pi_2(t)$ on $t$ at the event points, over the $50$
double-checked exponents (numbered $0,\dots,N-1$: the smallest exponent is the
origin, not an event), gives a slope of $2.6391$ and an intercept of $-2.9627$
with an OLS standard error of $0.4469$---apparently significant, and suggestive
of a secondary term in~\eqref{eq:LPW}. Nothing of the sort.

Two defects compound. The residuals of the cumulative count are strongly
autocorrelated, here $\rho_1=0.825$; the Newey--West covariance estimator with
eight lags raises the standard error to $0.7315$. More importantly, the
regression uses only the event points and thereby discards the empty interval
between the last exponent and the search frontier. The correct null model holds
the interval $[t_0,T]$ fixed, where $T=\ln X_{\mathrm{ver}}$, and conditions on
exactly $N-1=49$ events strictly to the right of $t_0$, which by
Proposition~\ref{prop:gamma} means $49$ independent uniform points on $(t_0,T]$.
Over $20\,000$ such replications the intercept has median $-1.498$ and a
$5$--$95\%$ range of $[-5.99,\,2.20]$, giving $p=0.562$ for the observed value.
The secondary term is an artifact of conditioning.

\subsection{Three consequences of one and the same conditioning}

\paragraph{(i) The bootstrap does not help.} The natural reaction to a
suspiciously small OLS standard error is to bootstrap over the event points. On
the base-2 data, a bootstrap over $20\,000$ resamples gives median $-2.949$,
standard deviation $0.387$ and a $95\%$ interval $[-3.822,\,-2.300]$ lying
entirely below zero; the proportion of negative values is $100\%$. This is not a
confirmation but a repetition of the same error: resampling the observed points
reproduces the variability within the realized configuration and cannot recover
the variability of the process that produced that configuration. The correct null
distribution is about six and a half times wider than the bootstrap one.

\paragraph{(ii) Spurious underdispersion.} Comparing the observed
$\max_i|r_i|=6.347$, where $r_i=i-\hat\lambda(t_i-t_0)$, with a simulation of an
\emph{unconditioned} Poisson process suggests the conclusion that the counting
function is ``smoother than Poisson''. Against the correct null
distribution---conditioning on $N-1$ events and estimating $\lambda$ from the
last of them---the median is $5.134$ and the two-sided significance level is
$0.539$. The observation is entirely ordinary.

\paragraph{(iii) The ``mid-range deficit''.} Estimating $\lambda$ from the last
event forces the residual $r$ to vanish at both ends ($r(t_0)=r(t_N)=0$ under
the correct numbering), turning its trajectory into an analogue of a Brownian
bridge, for which a negative excursion in the middle of the range arises of its
own accord. The observed mean excursion over the middle half of the range is
$-3.13$ against a simulation median of $-0.001$ and $p=0.280$; the proportion of
replications with a negative mean excursion is $0.500$, so there is no
systematic deficit. The trajectory of the residual process itself, together with
the simulation band, is shown in Figure~\ref{fig:bridge}.

\begin{figure}[t]
\centering
\includegraphics[width=.78\textwidth]{fig_resid_bridge.pdf}
\caption{The residual process $r(t)$ for base $b=2$ together with the
$5$--$95\%$ quantile band of $20\,000$ simulations and the simulation median.
The trajectory stays inside the band.}
\label{fig:bridge}
\end{figure}

\subsection{Base 2 at a known frontier}

Base~2 is the only one whose frontier is published, and hence the only one for
which a scheme~B analysis is possible with no assumptions at all. With
$T=\ln X_{\mathrm{ver}}$ we obtain $\hat\lambda=(N-1)/(T-t_0)=2.796$ against
$\eg/\ln2=2.570$, that is, $\hat\kappa_2=1.938$ with an exact $95\%$ Garwood
interval $[1.434,\,2.562]$ and a Poisson $p$-value of $0.59$ (doubled tail;
under the mid-$p$ convention, $0.54$---the convention must be stated here, since
the difference is appreciable). The conditional estimate on \emph{the same slice
of data}---the same $50$ double-checked exponents---gives $\tilde\kappa_2=1.905$;
the difference is small because the last exponent of that slice,
$77\,232\,917$, happens to lie very close to the frontier $81\,648\,221$. (For
comparison, the conditional estimate over all $52$ known exponents, reported in
Table~\ref{tab:bases}, is $1.921$; it must not be compared with $\hat\kappa_2$,
being a different sample.) This is a coincidence of one realization, not a
property of the method: the bias of \S\ref{sec:scheme} concerns the average over
the position of the last event relative to the frontier.

\subsection{The theoretical secondary term has the opposite sign}

There is an independent argument against a substantive reading of the intercept
$C_1$. If the secondary term in~\eqref{eq:LPW} originated in the Mertens
constant $\mathfrak M=0.26150\ldots$, its natural size would be
$(\eg/\ln2)\mathfrak M=+0.672$: the sign is positive, whereas what is observed is
$-2.963$. The nearest theoretical candidates not merely fail to reproduce the
observed magnitude---they predict the opposite sign. Independently of the
foregoing, this confirms that $C_1$ is a parameter of the fit and not of the
structure.

\section{Structural corrections}\label{sec:struct}

Consider the more general parametrization of the intensity
$\lambda_b=(\kappa/\ln b)\,C_b$ and four models for the factor:
\begin{align*}
\text{A:}&\quad C_b=1;\\
\text{B:}&\quad C_b=\prod_{p\mid b-1,\ p\ne2}\Bigl(\tfrac{p}{p-1}\Bigr)^{\alpha}
                 \prod_{p\mid b+1,\ p\ne2}\Bigl(\tfrac{p}{p-2}\Bigr)^{\beta};\\
\text{C:}&\quad C_b=\prod_{p\mid b^2-1}\Bigl(\tfrac{p}{p-1}\Bigr)^{\delta};\\
\text{D:}&\quad C_b=\exp\bigl(\theta_1\omega(b-1)+\theta_2\omega(b+1)+\theta_3\ln b\bigr).
\end{align*}
The exponent in model~C is written $\delta$ rather than $\gamma$ so as not to
confuse it with the Euler constant. The results are in Table~\ref{tab:struct}.
No parametrization gives a significant improvement; the BIC-based Bayes factors
overwhelmingly support the base model.

\begin{table}[t]
\centering
\caption{Comparison of structural models. The likelihood is the conditional
one~\eqref{eq:loglik} with factor $C_b$; LR and $p$ are relative to model~A;
BIC is computed with $n=M=218$.}
\label{tab:struct}
\small
\begin{tabular}{lrrrrrrr}
\toprule
Model & parameters & $\ln L$ & LR & $p$ & AIC & BIC & $P(\text{model}\mid\text{data})$\\
\midrule
A & 1 & $-176.129$ & --- & --- & 354.26 & 357.64 & 0.893\\
B & 3 & $-175.494$ & 1.27 & 0.530 & 356.99 & 367.14 & 0.008\\
C & 2 & $-175.647$ & 0.96 & 0.326 & 355.29 & 362.06 & 0.098\\
D & 4 & $-174.925$ & 2.41 & 0.492 & 357.85 & 371.39 & 0.001\\
\bottomrule
\end{tabular}
\end{table}

\begin{remark}[Sensitivity of the BIC]
The choice of effective sample size for the BIC is not obvious: the
likelihood~\eqref{eq:loglik} has $B=15$ independent summands, but the
information grows as $M=218$. With $n=15$ the posterior probabilities are
$0.624$ / $0.079$ / $0.261$ / $0.036$; with $n=218$ they are
$0.893$ / $0.008$ / $0.098$ / $0.001$. The conclusion favors model~A under
either choice, but with $n=218$ it is substantially stronger. We report $n=M$ as
matching the number of independent events.
\end{remark}

\paragraph{A fixed arithmetic correction.} We considered separately a model with
a \emph{fixed} (not estimated) factor
$C^{\mathrm{fix}}_b=\prod_{p\mid b-1,p\ne2}\frac{p}{p-1}\prod_{p\mid b+1,p\ne2}\frac{p}{p-2}$,
for which $\overline{C_b}=2.46$. It gives $\hat\kappa=0.815$. We stress that
comparing this value with $\eg$ is meaningless: in the reparametrized model
$\kappa$ has a different scale, and such a ``goodness-of-fit test'' would only
be testing whether the mean of $C_b$ equals one. Since models A and
$C^{\mathrm{fix}}$ are non-nested and have the same number of free parameters,
the correct comparison is directly by likelihood: $\ln L=-193.55$ against
$-176.13$, a difference of $-17.42$. The fixed correction is decisively
rejected.

\paragraph{Logarithmic corrections.} The model
$\lambda_b(t)=(\kappa/\ln b)(1+c_1/t+c_2/t^2)$ is estimated by the full
point-process likelihood $\sum_i\ln\lambda(t_i)-\int\lambda$ (the likelihood of
the count alone is inadmissible here: it discards the positions of the events,
which are the only thing that determines $c_1,c_2$). We obtain
$\hat\kappa=2.049$, $\hat c_1=+0.108$, $\hat c_2=-0.826$, $\mathrm{LR}=3.53$ on
two degrees of freedom, $p=0.171$. Not significant.

\section{Conclusions}

\begin{enumerate}
\item The estimators of the normalized frequency of generalized repunit primes
that are in common use are \emph{not} biased by the fitting scheme: both the
endpoint form and the OLS slope are exactly unbiased. Only the reciprocal is
biased, and that is a consequence of convexity. The common explanation of ``slow
convergence from below'' in terms of censoring does not hold.

\item The genuine defect of current practice is that different forms of the
estimator disagree by up to $29.5\%$, and that point values are published
without confidence intervals even though the OLS standard errors in this problem
are understated by a factor of about $1.6$. The classification of bases into
``lucky'' and ``unlucky'' flips when the form of the estimator is changed on
unchanged data.

\item Pooling $B=15$ bases yields $M=218$ events and the exact pivot
$\kappa S\sim\Gam(M,1)$, but that pivot is calibrated to the observation scheme
``until the $N$-th event'', whereas the data arise from a fixed search frontier.
Under the actual scheme a nominal $5\%$ test rejects in $16.2\%$ of cases, and
the correct pooled estimator is $(M-B)/S$: the conditional likelihood costs one
event per base, not one event in total. Distinguishing such schemes goes back to
reliability theory; what is new here is the case of an undocumented frontier,
where the exposure cannot be taken from a protocol and has to be reconstructed
from the finds (Remark~\ref{rem:priority}).

\item This gives $\kappa=1.838$ with a $95\%$ interval $[1.61,\,2.13]$ and
$p=0.639$ against $\eg$. The reported ten-percent excess over LPW---including
our own in an earlier version of this work---is an artifact of the observation
scheme. Moving to scheme~B independently removes the systematic trend of the
pooled estimate under truncation in $n_{\min}$.

\item The choice between $(M-1)/S$ and $(M-B)/S$ is not settled by the data:
both estimators are functions of the same sufficient statistics and differ by
$7\%$. The practical consequence is the need to publish the search frontiers
$L_b$, as GIMPS does for $b=2$. Until that is done, any pooled analysis of these
sequences carries an irreducible uncertainty of order $7\%$---twice the residual
gap with LPW.

\item The LPW conjecture is not rejected by the available data, and
base-dependent corrections are not required. But both of these are statements
about the low sensitivity of the tests: the homogeneity test resolves only a
spread of $\kappa_b$ of order $\pm34\%$, and the pooled test only deviations of
the constant of at least $21\%$. Attaining a sensitivity of $10\%$ would require
the search frontiers to grow by a factor of about $3\cdot10^{20}$.
\end{enumerate}

\appendix
\section{Verification of the source data}\label{app:data}

All twenty sequences (fifteen of the main set and five from the check of the
selection limit) were cross-checked against the OEIS entries on 27 August 2026.
Since only the triples $(N_b,n_{\min},n_{\max})$ enter the pooled analysis,
Table~\ref{tab:verify} gives, for the main set, the number of known primes---from
which $N_b$ is obtained by subtracting one---and both extreme indices; the full
lists of indices are contained in the script \texttt{analysis.py}.

\begin{table}[h]
\centering
\caption{Cross-check against OEIS. The column ``last addition'' gives the
provenance of the largest known index.}
\label{tab:verify}
\small
\begin{tabular}{rlrrrl}
\toprule
$b$ & OEIS & primes & $n_{\min}$ & $n_{\max}$ & last addition\\
\midrule
2 & A000043 & 52 & 2 & 136\,279\,841 & GIMPS, Oct.\ 2024\\
3 & A028491 & 23 & 3 & 8\,530\,117 & P.~Bourdelais\\
5 & A004061 & 19 & 3 & 3\,300\,593 & P.~Bourdelais\\
6 & A004062 & 17 & 2 & 3\,360\,347 & P.~Bourdelais\\
7 & A004063 & 10 & 5 & 1\,264\,699 & P.~Bourdelais\\
10 & A004023 & 11 & 2 & 8\,177\,207 & Propper and Batalov, May 2021\\
11 & A005808 & 13 & 17 & 1\,868\,983 & P.~Bourdelais\\
12 & A004064 & 14 & 2 & 769\,543 & P.~Bourdelais, Dec.\ 2014\\
13 & A016054 & 13 & 5 & 1\,503\,503 & P.~Bourdelais\\
14 & A006032 & 11 & 3 & 1\,724\,417 & P.~Bourdelais, 13 Jul.\ 2026\\
15 & A006033 & 10 & 3 & 639\,833 & P.~Bourdelais, 22 Apr.\ 2019\\
17 & A006034 & 12 & 3 & 1\,990\,523 & P.~Bourdelais, 3 Aug.\ 2020\\
18 & A133857 & 8 & 2 & 1\,270\,141 & P.~Bourdelais\\
19 & A006035 & 11 & 19 & 209\,359 & P.~Bourdelais, 27 Aug.\ 2010\\
20 & A127995 & 9 & 3 & 984\,349 & P.~Bourdelais\\
\bottomrule
\end{tabular}
\end{table}

A separate warning to the reader who checks the data against secondary sources:
summary tables (the ``Repunit primes'' page in the OEIS Wiki, the article ``List
of repunit primes'' in Wikipedia) are updated with a delay and, at the time of
writing, lagged the OEIS entries themselves by one term for the bases
$b=14,15,17,18,20$. The OEIS entries should be treated as authoritative.

\subsection*{Independent recomputation of the prefixes}

Cross-checking against OEIS confirms that the listed indices give primes, but it
does not guarantee \emph{completeness}: a missing term would shift $N_b$, and
with it both pooled estimates. We closed this gap by direct recomputation. For
each of the twenty bases $2\le b\le26$ that are not perfect powers, an
independent implementation\footnote{A hybrid CPU/GPU searcher: segmented sieving
by small divisors, trial factoring in Montgomery arithmetic, PRP testing via GMP
and GWNUM/IBDWT. The source code and the run logs are in the same repository,
scripts \texttt{verify\_sequences.sh} and \texttt{compare\_verify.py}.} tested
\emph{all} indices $n<10^4$ consecutively, not only those listed in OEIS. Every
negative verdict was recomputed by a second, independent backend (a blanket
double-check, as in GIMPS) so as to rule out a lost find.

The result: $153$ terms confirmed, none spurious, none missing, no disagreement
between backends---across all twenty bases. Coverage is uneven across bases
(from one term at $b=18$ to twenty-two at $b=2$): the main set $b\le20$ accounts
for $128$ confirmed indices out of $233$ known. Above $10^4$ a complete
recomputation is computationally out of reach, and there we still rely on OEIS.

We note that this recomputation was not a formality: on the first run the
independent implementation declared $R_{3181}(23)$ composite, disagreeing with
A204940. Analysis showed an error in the implementation itself---a spike of FFT
rounding error on the last iteration, after the last periodic check---and not in
the OEIS data. After the fix no disagreements remained. We report this episode as
an illustration of the thesis of \S\ref{sec:scheme}: a false ``composite'' is
invisible without recomputation, and that is exactly why frontiers and
verification protocols should be published.

\section*{Declarations}

\subsection*{Acknowledgments}
The author thanks the Great Internet Mersenne Prime Search (GIMPS) project for
providing the Mersenne prime data through the open PrimeNet report
(\url{https://www.mersenne.org/}); the Mersenne exponents were retrieved on
27 August 2026. The author is grateful to the OEIS contributors who have
maintained the generalized repunit prime sequences for many years.

\subsection*{Funding}
This research received no external funding.

\subsection*{Competing interests}
The author declares no competing interests.

\subsection*{Data and code availability}
All sequences used are listed in Table~\ref{tab:data}. The repository
\url{https://github.com/Wilps93/repunit-hunt}, directory \texttt{paper/},
contains everything needed to reproduce the paper; the file tree and the order
of execution are described in \texttt{paper/README.md}.

\emph{Data.} The script \texttt{fetch\_data.sh} downloads the OEIS b-files for
all fifteen sequences of the main set, and \texttt{fetch\_extra.sh} does the
same for the five bases $21\le b\le26$ used in the check of the selection limit
(\S\ref{sec:bases}). The downloaded files are included in the repository, so the
analysis reproduces without network access.

\emph{Analysis.} A single self-contained script \texttt{analysis.py} (numpy,
scipy, matplotlib); every number in this paper is tagged in its output
(\texttt{results.txt}) with the corresponding table or section and is reproduced
by one run. The random number generator is initialized with a fixed seed, so all
Monte Carlo quantities reproduce bit for bit. The script also produces all four
figures. A full run takes a few minutes.

\emph{Independent recomputation of the sequences.}
\texttt{verify\_sequences.sh} runs the searcher over all indices $n<10^4$ for
each base, and \texttt{compare\_verify.py} compares the result with the OEIS
b-files, returning a non-zero exit code on any disagreement
(Appendix~\ref{app:data}).

\emph{A caveat on the Monte Carlo quantities.} The endpoints of
interval~\eqref{eq:resB} depend on the inversion procedure, and
Tables~\ref{tab:mc}, \ref{tab:frontier} and~\ref{tab:calib} depend on the rules
for placing the frontier. Both are described in the text and implemented in
\texttt{analysis.py}; every value reported in this paper is exactly the output
of that script under the seed fixed in it. Simulation accuracy: about $\pm0.003$
for the significance level in Table~\ref{tab:calib}, and about $\pm0.3\%$ for
the means in Table~\ref{tab:frontier}.

\begin{thebibliography}{99}

\bibitem{wagstaff1983} S.~S. Wagstaff, Jr., \emph{Divisors of Mersenne
numbers}, Math. Comp. \textbf{40} (1983), no.~161, 385--397.
\href{https://doi.org/10.1090/S0025-5718-1983-0679454-X}{DOI:
10.1090/S0025-5718-1983-0679454-X}.

\bibitem{pomerance1981} C.~Pomerance, \emph{On the distribution of
pseudoprimes}, Math. Comp. \textbf{37} (1981), no.~156, 587--593.
\href{https://doi.org/10.1090/S0025-5718-1981-0628717-0}{DOI:
10.1090/S0025-5718-1981-0628717-0}.

\bibitem{t5k} \emph{Heuristics: deriving the Wagstaff Mersenne conjecture}, The
PrimePages: prime number research and records.
\url{https://t5k.org/mersenne/heuristic.html} (accessed 27 August 2026).

\bibitem{dubner1993} H.~Dubner, \emph{Generalized repunit primes}, Math. Comp.
\textbf{61} (1993), no.~204, 927--930.
\href{https://doi.org/10.1090/S0025-5718-1993-1185243-9}{DOI:
10.1090/S0025-5718-1993-1185243-9}.

\bibitem{bourdelais2009} P.~Bourdelais, \emph{A generalized repunit
conjecture}, NMBRTHRY mailing list, 25 June 2009.
\url{https://listserv.nodak.edu/cgi-bin/wa.exe?A2=NMBRTHRY;417ab0d6.0906}
(accessed 25 August 2026).

\bibitem{oeis} OEIS Foundation Inc., \emph{The On-Line Encyclopedia of Integer
Sequences}. \url{https://oeis.org} (accessed 27 August 2026).

\bibitem{wagstaff2025} S.~S. Wagstaff, Jr., \emph{Two conjectures about Mersenne
primes}, J. Integer Seq. \textbf{28} (2025), article 25.7.2.
\url{https://cs.uwaterloo.ca/journals/JIS/VOL28/Wagstaff/wag3.html}
(accessed 25 August 2026).

\bibitem{dominguez2026} J.~Dominguez, \emph{Arithmetic bias in the distribution
of Mersenne prime exponents and the divisor structure of $p-1$}, preprint,
arXiv:2603.08994 (2026). \url{https://arxiv.org/abs/2603.08994}
(accessed 27 August 2026).

\bibitem{jha2026} A.~Jha, \emph{The Poisson tail conjecture for primes in short
intervals}, preprint, arXiv:2605.23014 (2026).
\url{https://arxiv.org/abs/2605.23014} (accessed 27 August 2026).

\bibitem{gimps} Great Internet Mersenne Prime Search (GIMPS), \emph{PrimeNet};
Prime95/MPrime software, version 30.8. \url{https://www.mersenne.org/}
(accessed 27 August 2026; the \emph{Milestones} report was read at
27 August 2026, 04:30~UTC).

\bibitem{epstein1954trunc} B.~Epstein, \emph{Truncated life tests in the
exponential case}, Ann. Math. Statist. \textbf{25} (1954), no.~3, 555--564.
\href{https://doi.org/10.1214/aoms/1177728723}{DOI: 10.1214/aoms/1177728723}.

\bibitem{epsteinsobel1954} B.~Epstein and M.~Sobel, \emph{Some theorems
relevant to life testing from an exponential distribution}, Ann. Math. Statist.
\textbf{25} (1954), no.~2, 373--381.
\href{https://doi.org/10.1214/aoms/1177728793}{DOI: 10.1214/aoms/1177728793}.

\bibitem{bartholomew1963} D.~J. Bartholomew, \emph{The sampling distribution of
an estimate arising in life testing}, Technometrics \textbf{5} (1963), no.~3,
361--374.

\bibitem{basu1964} A.~P. Basu, \emph{Estimates of reliability for some
distributions useful in life testing}, Technometrics \textbf{6} (1964), no.~2,
215--219.

\bibitem{kivela2015} M.~Kivel\"{a} and M.~A. Porter, \emph{Estimating interevent
time distributions from finite observation periods in communication networks},
Phys. Rev. E \textbf{92} (2015), no.~5, 052813.
\href{https://doi.org/10.1103/PhysRevE.92.052813}{DOI:
10.1103/PhysRevE.92.052813}.

\end{thebibliography}

\end{document}
